\documentclass[11pt]{article}
\usepackage{amsmath,amssymb,amsthm,mathtools}
\usepackage[margin=1in]{geometry}
\usepackage{hyperref}
\usepackage{enumitem}

\newcommand{\Ext}{\operatorname{Ext}}
\newcommand{\Hom}{\operatorname{Hom}}
\newcommand{\im}{\operatorname{im}}
\newcommand{\coker}{\operatorname{coker}}
\newcommand{\id}{\mathrm{id}}
\newcommand{\A}{\mathcal{A}}
\newcommand{\F}{\mathbb{F}}
\newcommand{\kk}{\Bbbk}
\newcommand{\tcofree}[1]{\mathcal{F}(#1)}
\theoremstyle{definition}
\newtheorem{remark}{Remark}

\title{Yoneda products for the $C$-motivic minimal resolution\\
\large An implementation note for \texttt{yoneda.cpp}}
\author{}
\date{\today}

\begin{document}
\maketitle

\begin{abstract}
This note explains the mathematics behind \texttt{yoneda.cpp}, a tool that computes
Yoneda (composition) products on the $\Ext$ groups produced by Guozhen Wang's
\texttt{MinimalResolution} infrastructure for the $C$-motivic Steenrod algebra.
The existing code already computed multiplication by the filtration-$1$ classes
$h_n$ (this is \texttt{mot\_mult}); we generalize this to multiplication by an
\emph{arbitrary} class via chain-map lifting. I assume comfort with homological
algebra (resolutions, $\Ext$, Yoneda products) but treat the motivic apparatus as
a black box, calling out only the two features that actually change the algorithm:
the ground ring is $\F_2[\tau]$ rather than a field, and the relevant modules can
have $\tau$-torsion.
\end{abstract}

\section{The homological setup, with the motivic part boxed off}

\paragraph{What is being resolved.}
Fix the ground ring $\kk=\F_2[\tau]$, a graded polynomial ring on one generator
$\tau$. (Black box: $\kk$ is the mod-$2$ motivic cohomology of a point; $\tau$ is a
non-nilpotent element of homological degree $0$. Setting $\tau=1$ recovers the
classical mod-$2$ story, and ``$\tau$-torsion'' is the genuinely new phenomenon.)
There is a graded Hopf algebroid $(\kk,\A)$ -- the $C$-motivic dual Steenrod
algebra -- and we work in the category of $\A$-comodules. The object of interest is
\[
  \Ext_{\A}^{s,*}(\kk,\kk),
\]
the $E_2$-page of the $C$-motivic Adams spectral sequence for the sphere. Each
group is a graded $\kk=\F_2[\tau]$-module.

\paragraph{Cofree comodules and the resolution.}
For a graded $\kk$-module $V$ there is a \emph{cofree} (relatively injective)
comodule $\tcofree{V}=\A\,\square\,V$, characterized by the adjunction
\begin{equation}\label{eq:adjunction}
  \Hom_{\A}\bigl(M,\tcofree{V}\bigr)\;\cong\;\Hom_{\kk}(M,V),
\end{equation}
\emph{natural in the comodule $M$}: a comodule map into $\tcofree V$ is the same
thing as a $\kk$-linear map into the ``cogenerators'' $V$. Explicitly, a
$\kk$-linear $\psi\colon M\to V$ corresponds to the comodule map
\begin{equation}\label{eq:adjoint}
  \widehat\psi(m)=(\id_\A\otimes\psi)\,\bigl(\psi_M(m)\bigr)\in \A\square V=\tcofree V,
\end{equation}
where $\psi_M\colon M\to \A\otimes M$ is the coaction. Formula~\eqref{eq:adjoint} is
the computational workhorse below.

\texttt{MinimalResolution} builds a \emph{minimal} resolution of the trivial
comodule $\kk$ by cofree comodules,
\begin{equation}\label{eq:resolution}
  0\to \kk \xrightarrow{\ } F_0 \xrightarrow{d_0} F_1 \xrightarrow{d_1} F_2 \to\cdots,
  \qquad F_s=\tcofree{V_s},
\end{equation}
i.e.\ an injective resolution in comodules. Write $G_s$ for the cogenerators of
$F_s$ (a finite free $\kk$-module; in the code these are
\texttt{F.generators}). Because the resolution is minimal, the induced differential
on cogenerators vanishes \emph{modulo $\tau$}, and one has, up to the
$\tau$-Bockstein bookkeeping discussed in \S\ref{sec:naming},
\[
  \Ext_{\A}^{s}(\kk,\kk)\;\cong\;G_s
  \qquad(\text{a basis is the cogenerators of }F_s).
\]
A class is therefore named by a pair $\{s\text{-}a\}$: homological filtration $s$
and the index $a$ of a cogenerator of $F_s$. For example $h_n=\{1\text{-}n\}$.

\paragraph{How the resolution is stored.}
Each step is realized by a short exact sequence of comodules
\begin{equation}\label{eq:ses}
  0\to X_s \xrightarrow{\ \iota_s\ } F_s \xrightarrow{\ q_s\ } X_{s+1}\to 0,
  \qquad X_0=\kk,
\end{equation}
with the resolution differential $d_s=\iota_{s+1}\circ q_s\colon F_s\to F_{s+1}$.
The comodules $X_s$ are the (co)syzygies. On disk one finds, for each $s$: the
cofree comodule $F_s$ (file \texttt{mot\_gens}$s$) and the two $\kk$-linear maps
$\iota_s,q_s$ (file \texttt{mot\_maps}$s$), all as sparse matrices over $\F_2[\tau]$.
Every map in sight is an honest $\F_2[\tau]$-linear map between finitely generated
free $\F_2[\tau]$-modules; the comodule structure enters only through the
coactions, via \eqref{eq:adjoint}.

\section{Yoneda products via chain-map lifting}

\paragraph{The classical recipe.}
A class $\beta\in\Ext^{s'}(\kk,\kk)$ is represented by a map $\kk\to\kk[s']$ in the
derived category, equivalently (using \eqref{eq:resolution} as an injective
resolution of the target) by a \emph{chain map}
\[
  f_\bullet\colon F_\bullet \longrightarrow F_{\bullet+s'},
  \qquad d_{i+s'}\,f_i = f_{i+1}\,d_i,
\]
that lifts $\beta$ in degree $0$. The Yoneda product $\alpha\cdot\beta$, for
$\alpha\in\Ext^{s}$, is obtained by applying $f_s$ to a cocycle representing
$\alpha$ and reading the result in $\Ext^{s+s'}$. Two chain maps lifting $\beta$
differ by a homotopy, so the induced map on cohomology -- the product -- is
well defined. This is exactly the ``simplest approach'': build $f_\bullet$ by
lifting, one filtration at a time.

\paragraph{The base of the lift.}
$F_0=\tcofree{\kk}$ is cofree on a single cogenerator $u$ in degree $0$. We set
$f_0$ to be the comodule map sending $u$ to the cogenerator $b$ of $F_{s'}$ that
\emph{is} $\beta=\{s'\text{-}b\}$. Via \eqref{eq:adjunction} this is encoded by the
$\kk$-linear map $\psi_0\colon F_0\to G_{s'}$ with $\psi_0(u)=e_b$ and $\psi_0=0$ on
the rest of a $\kk$-basis of $F_0$; then $f_0=\widehat{\psi_0}$ by
\eqref{eq:adjoint}.

\paragraph{The inductive step.}
Assume $f_i\colon F_i\to F_{i+s'}$ is built. We must produce
$f_{i+1}\colon F_{i+1}\to F_{i+s'+1}$ with $f_{i+1}d_i=d_{i+s'}f_i$. The key points:

\begin{enumerate}[leftmargin=*]
\item \textbf{$f_{i+1}$ is forced on the image of $d_i$.}
  Since $d_i=\iota_{i+1}q_i$ and $\ker d_i=\ker q_i=\im\iota_i$, the chain-map
  equation determines $f_{i+1}$ on $\im d_i=\im\iota_{i+1}$ by
  \begin{equation}\label{eq:forced}
    f_{i+1}\bigl(d_i(e)\bigr)=d_{i+s'}\bigl(f_i(e)\bigr),
    \qquad e\in F_i,
  \end{equation}
  and this is consistent precisely because $\beta$ is a cocycle (the obstruction
  $d_{i+s'}f_i|_{\im\iota_i}$ vanishes by the previous step).
\item \textbf{Cogenerators live in that image.} The cogenerators of $F_{i+1}$ are
  the leading terms of $\iota_{i+1}$, hence lie in $\im\iota_{i+1}=\im d_i$. So
  \eqref{eq:forced} computes $f_{i+1}$ \emph{on cogenerators}, which by
  \eqref{eq:adjoint} is all we need to reconstruct $f_{i+1}$ as a comodule map.
\end{enumerate}
Concretely, to compute $f_{i+1}$ on the cogenerator $w$ (a basis vector of
$F_{i+1}$):
\begin{equation}\label{eq:step}
  u:=\iota_{i+1}^{-1}(w)\in X_{i+1},\quad
  e:=q_i^{-1}(u)\in F_i,\quad
  f_{i+1}(w):=d_{i+s'}\bigl(f_i(e)\bigr).
\end{equation}
Here $\iota_{i+1}^{-1}$ means the unique preimage under the \emph{injection}
$\iota_{i+1}$ (it exists because $w\in\im\iota_{i+1}$), and $q_i^{-1}$ means
\emph{any} preimage under the surjection $q_i$; then
$d_i(e)=\iota_{i+1}(q_i(e))=\iota_{i+1}(u)=w$, so \eqref{eq:step} is just
\eqref{eq:forced}. Finally we package $\{w\mapsto f_{i+1}(w)\}$ back into the full
comodule map via the adjoint \eqref{eq:adjoint}; and $f_i(e)$ for the general
element $e$ in \eqref{eq:step} is itself evaluated by \eqref{eq:adjoint} applied to
$f_i=\widehat{\psi_i}$.

\paragraph{Reading off the product.}
This yields, on cogenerators, a $\kk$-linear ``multiplication matrix''
$M_s\colon G_s\to G_{s+s'}$, $M_s=\bigl(\text{project to cogenerators}\bigr)\circ
f_s$. For $\alpha=\{s\text{-}a\}$,
\[
  \alpha\cdot\beta \;=\; M_s\bigl(\text{cocycle of }\alpha\bigr)\ \in\ \Ext^{s+s'},
\]
re-expressed in the chosen basis of $\Ext^{s+s'}$ (\S\ref{sec:naming}). Specializing
$\beta=h_n$ recovers exactly the construction in \texttt{mot\_mult} (there the
chain map is the elementary one ``multiply by the algebra element $h_n$, then apply
one differential''); our code reproduces those tables identically as a sanity check.

\section{The one motivic complication: $\tau$-torsion and $\tau$-aware solving}

Everything above is standard homological algebra. The single place where ``motivic''
intervenes is in computing the two preimages in \eqref{eq:step}, i.e.\ in inverting
\[
  \iota_{i+1}\ \text{on its image}\qquad\text{and}\qquad q_i\ \text{(a surjection).}
\]
Over a field this is ordinary Gaussian elimination. Over $\kk=\F_2[\tau]$ it is not,
and naive elimination fails.

\paragraph{Why a field-style splitting does not exist.}
$\F_2[\tau]$ is a PID, and a short exact sequence \eqref{eq:ses} of $\kk$-modules
splits iff $X_{s+1}=\coker\iota_s$ is free, i.e.\ \emph{$\tau$-torsion-free}.
$C$-motivic $\Ext$ is famous for having $\tau$-torsion, so the syzygies are in
general not free, and $\iota_s$ has \emph{no global retraction}
$F_s\to X_s$. (Empirically: in the $10$-cell test resolution, attempting a global
$\tau^0$-pivoted retraction already fails at step $1$, on the row whose $\iota_1$-image
is $\tau$-divisible.) This is exactly why the existing code never lifts chain maps
over $\F_2[\tau]$ -- and the conceptual obstacle the present work had to get past.

\paragraph{The resolution of the difficulty.}
We never need a global retraction. In \eqref{eq:step} we only ever
\begin{itemize}[leftmargin=*]
\item invert $\iota_{i+1}$ on an element $w$ \emph{known to lie in $\im\iota_{i+1}$}
  (a cogenerator); since $\iota_{i+1}$ is injective the preimage is unique and
  \emph{always exists}; and
\item invert $q_i$ on an arbitrary target element; since $q_i$ is surjective a
  preimage \emph{always exists}.
\end{itemize}
Existence is guaranteed by exactness of \eqref{eq:resolution}; there is no torsion
obstruction to \emph{solving} these particular systems. What is required is an
elimination that works over the PID $\F_2[\tau]$.

\paragraph{$\tau$-aware (Hermite) elimination.}
The implementation builds, for each matrix $\iota_i$ and $q_i$, a row-echelon
decomposition over $\F_2[\tau]$ in which the pivot at each position is kept with the
\emph{minimal} $\tau$-power seen so far (a one-prime Hermite normal form; the only
relevant prime is $\tau$). Reducing a target $w$ against these pivots either
certifies $w\notin\im$ (a residual survives) or returns a preimage, with the
elimination introducing $\tau^{k}$ factors $(k\ge 0)$ exactly when the system forces
them. This is the class \texttt{TauEchelon} in the code; the same routine serves both
$\iota^{-1}|_{\im}$ and the section of $q$.

\paragraph{Where the $\tau$'s in products come from.}
The $\tau$-powers produced by this elimination are precisely the source of
$\tau$-multiples in Yoneda products. The cleanest example: classically
$h_0^2h_2=h_1^3$, while our computation returns
\[
  h_1^3=\{3\text{-}1\},\qquad h_0^2h_2=\tau\cdot\{3\text{-}1\},
\]
i.e.\ the $C$-motivic refinement $\boxed{\,h_0^2h_2=\tau\,h_1^3\,}$, with the $\tau$
appearing exactly as the Hermite elimination's $\tau^1$ factor.

\section{Naming the answer: the $\tau$-Bockstein basis}\label{sec:naming}

One subtlety remains in how $\Ext$ is presented. Because the minimal differential is
$\tau$-divisible, the honest $\F_2[\tau]$-module $\Ext^{s}$ is not simply ``the
cogenerators of $F_s$'': one must run the $\tau$-Bockstein spectral sequence to find
the correct cycle representatives and the $\tau$-towers. This is already implemented
(\texttt{make\_table}/\texttt{make\_cycle\_tables} in \texttt{tao\_bockstein.cpp}),
producing for each filtration a distinguished set of cycle representatives
(\texttt{cyc\_table}). Our product routine names its output in this basis, via the
existing \texttt{find\_cycle}, so that the result is directly comparable with the
$h_n$-multiplication tables and with the literature. The $\tau$-exponent printed as
\texttt{t\^{}$k$} in front of a generator is the $\tau$-power attached by this naming.

\section{Summary of the algorithm and verification}

\paragraph{Algorithm (file \texttt{yoneda.cpp}).}
\begin{enumerate}[leftmargin=*,itemsep=1pt]
\item Load the resolution $F_\bullet$ and the maps $\iota_s,q_s$; build $\tau$-aware
  solvers \texttt{invInj}$_s$ (invert $\iota_s$ on its image) and \texttt{secQut}$_s$
  (section of $q_s$). A self-check verifies
  $\iota_s^{-1}\iota_s=\id$ and $q_s\,q_s^{-1}=\id$ on the relevant spaces.
\item Build the $\tau$-Bockstein cycle representatives (reusing
  \texttt{make\_table}).
\item For the chosen $\beta=\{s'\text{-}b\}$, lift the chain map $f_\bullet$ by
  \eqref{eq:step}, using \eqref{eq:adjoint} to evaluate and to reassemble comodule
  maps. Record the cogenerator-level matrices $M_s$.
\item Output either a single product $\alpha\cdot\beta$, or the full table
  $\beta\cdot(-)$ over all generators, named in the $\tau$-Bockstein basis.
\end{enumerate}
A technical point in step~3: when a product's internal degree exceeds the resolution
range, \eqref{eq:adjoint} would place terms outside the relevant cofree summand;
these are degree-overflow artifacts and are clipped to the summand, which is correct
within range (nothing valid is dropped) and matches \texttt{mot\_mult}'s degree
guard.

\paragraph{Verification.}
On a test resolution (\texttt{md=10}, length $6$):
\begin{itemize}[leftmargin=*,itemsep=1pt]
\item The solver self-check passes at every filtration.
\item Taking $\beta=h_n$ reproduces \texttt{mot\_mult}'s tables
  $\{n=0,1,2,3\}$ \emph{exactly}, including $\tau$-multiples such as
  $h_0\cdot\{2\text{-}2\}=\tau\{3\text{-}1\}$ and $h_0h_1=0$.
\item Genuinely higher-filtration (input) products agree with known relations:
  $h_1^2=\{2\text{-}1\}$, $h_1^3=\{3\text{-}1\}$, and
  $h_0^2h_2=\tau\{3\text{-}1\}=\tau h_1^3$.
\item Single-pair queries agree with the corresponding entries of the full tables.
\end{itemize}

\paragraph{Usage.}
After the bootstrap \texttt{e2p $md$;\ motTab $md$;\ mr\_ex $md$ $L_{\mathrm{ex}}$;\
mr\_mot $md$ $L$} (note: \texttt{motTab} is required), run
\begin{center}
\begin{tabular}{ll}
\texttt{yoneda $md$ $L$} & solver self-check\\
\texttt{yoneda $md$ $L$ $b_s$ $b_b$} & full table for $\beta=\{b_s\text{-}b_b\}$\\
\texttt{yoneda $md$ $L$ $a_s$ $a_a$ $b_s$ $b_b$} & single product $\{a_s\text{-}a_a\}\cdot\{b_s\text{-}b_b\}$\\
\end{tabular}
\end{center}
where $md$ is half the maximal topological degree and $L$ the resolution length.

\begin{remark}[scope]
Only Yoneda products are implemented. Massey products are a natural sequel: the
$\tau$-aware solving developed here already produces the null-homotopies of
vanishing products that a Massey-product implementation would consume; what remains
is the indeterminacy bookkeeping.
\end{remark}

\end{document}
